\subsection{最後の単冪Jordan型\texorpdfstring{$(2,2)$}{(2,2)}と単一固有値のCY性}
\label{subsec:cubic-jordan22-unipotent-vertex}
\begin{theorem}\label{thm:cubic-jordan22-unipotent-empty}
強い内在的H3の二次包絡の中山作用は、単一固有値1のJordan型 $(2,2)$ を持たない。
従って中山作用の固有値が一種類しかないH3包絡は、必ずCYである。
\end{theorem}
\begin{proof}
標準ねじり解除を $B=C^H$、$C$ CY、$H=\exp N$ と書き、
$Nx_1=x_0$, $Nx_3=x_2$, $Nx_0=Nx_2=0$ とする。
$U=\langle x_0,x_2\rangle$、$V=\langle x_1,x_3\rangle$ とし、$V$ の個数をweightとする。
元反巡回topの $N_{\rm tot}$ 核は26次元であり、weight0,1,2の次元はそれぞれ4,12,10である。
次の語の係数を順に $a_0,\ldots,a_{25}$ とする：
\[
\begin{gathered}
0001,0002,0003,0011,0012,0013,0021,0022,0023,0031,0032,0033,0122,\\
0123,0202,0203,0213,0222,0223,0232,0233,0303,0322,0323,2223,2233.
\end{gathered}
\]
middleの $U^2$ 行・$V^2$ 列blockは
\[
 Q=\begin{pmatrix}
 a_3&a_5&a_9&a_{11}\\
 a_5&a_{16}&a_{11}+a_{13}&a_{20}\\
 a_9&a_{11}+a_{13}&a_{16}+a_{21}&a_{20}+a_{23}\\
 a_{11}&a_{20}&a_{20}+a_{23}&a_{25}
 \end{pmatrix}.
\]
その列は純 $U^2$ の実関係である。
二次元の二次関係空間は代数閉体上でdecomposable tensorを含むので、整域性から
$\operatorname{rank}Q\le1$。
$Q=0$ なら全weight2部分が零となる。このときmiddleはweight0とweight1で
$\left(\begin{smallmatrix}P&M\\M^t&0\end{smallmatrix}\right)$ の形となり、
rank6から $r=\operatorname{rank}M\le3$。
全関係の純 $U^2$ 部分は次元 $6-r\ge3$ となり、再び整域性に反する。
従って $Q$ は非零rank1である。

対称性と代数閉性により $Q=qq^t$ と書ける。
$q$ 自体も実二次関係なので、その $2\times2$ 係数行列は可逆である。
$U$ の基底変換と $V$ の同じ基底変換は $N$ を保つ。
二生成元AS2の二次標準形により、$q$ は
\[
 q=(0,1,-p,0),\ p\ne0,\qquad\text{又は}\qquad q=(1,1,-1,0)
\]
にできる。topの尺度も正規化した。
前者の最高mixed blockの一つの五次minorは
$-p^2(p-1)^2(p+1)^2$ である。
純 $U^2/V^2$ のrank1対合が既にrank2を使うので、middle全rank6はmixed rankが高々4であることを要求する。
従って $p=1,-1$ だけが残る。
以上で最高weightの全層は、可換・反可換・Jordanの三表示で漏れなく覆われる。

各表示でmiddleの行列番号 $(2,7,1,3,6,11)$ の六次principal minorは $-1$ である
（番号 $4i+j$ は語 $ij$ を表す）。
従って全低weight係数のまま元六関係を構成できる。
最高関係の六逆順pair $(20,10,12,30,32,31)$ の係数行列も単元であり、
$U$ 内と $V$ 内はそれぞれ同じAS2関係、四つの $VU$ 対は $UV$ の一次結合へ書き換えられる。
まず $V/U$ の並べ替え、次に各二元の並べ替えを行えば、
$x_0^ix_2^jx_1^kx_3^l$ が張る。
その個数はAS4 Hilbert値と一致するので、実際の $C$ のweight随伴環は
この最高関係の代数と全次数で一致する。
特に退化した低weight係数でもこの議論は変わらない。

この三つの最高代数の三次商は20次元で、誘導された $N$ の核は各々6次元である。
その全核のweightは高々1となる。
これは元の最高六関係を用いた $20\times20$ 有理行列の核として厳密に検証する。
従って輸送された正常三次 $N(F)=0$ の先頭もweightが高々1であり、
$J=(x_0,x_1,x_2)$ に対して $F\in J^2$ である。
三表示とも $C/J=k[x_3]$、$H(J)=J$、$H(x_3)\equiv x_3\pmod J$ なので、
$R=B/(F)$ に $P=\pi k[x_3]$ がある。

元の全低weight係数を残してpointの自己拡大を一次化する。
$x_3$ を恒等とし、他の非対角成分を $(f_0,f_1,f_2)$ とする。
全六二次関係は、三つの独立な漸化式に一致し、他の一次条件は無い。
$F\in J^2$ も追加条件を課さないので $\Ext^1(P,P)=k^3$ である。
$B=C^H$ へtwistした後のshift行列は次の形になる：
\[
\begin{array}{c|c|c}
\text{表示}&M&\det M\\\hline
\text{可換}&\begin{pmatrix}1&0&0\\\lambda&1&\rho\\0&0&1\end{pmatrix}&1\\
\text{反可換}&\begin{pmatrix}-1&0&0\\\lambda&-1&\rho\\0&0&1\end{pmatrix}&1\\
\text{Jordan}&\begin{pmatrix}1&0&0\\\lambda&1&\rho\\3&0&1\end{pmatrix}&1
\end{array}
\]
ここで可換では $\lambda=a_{15}-a_8+1$, $\rho=a_{18}+a_{19}$、
反可換では $\lambda=-a_{15}-a_8-1$, $\rho=-a_{18}-a_{19}$、
Jordanでは $\lambda=a_{15}+3a_{19}-a_8+1$, $\rho=a_{18}+a_{19}$。
三つの漸化式を解くminorも単元なので、これは全係数・全特殊値で成立する。
反可換の場合にshiftが単冪であるとは主張しない。
必要なのは行列式1であり、補題\ref{lem:cubic-odd-Serre-determinant}が
三次元の自己拡大とH3のSerre作用の両立を否定する。
これで $(2,2)$ も全排除した。

単一固有値は命題\ref{prop:h3-spectral-filter}より1である。
非自明単冪の四つのJordan型 $(4),(3,1),(2,2),(2,1,1)$ は本定理と前節までで全て排除された。
従って単一固有値の場合は $\mu_B=I$、すなわちCYである。
\end{proof}
公開再現器は\\
\texttt{verify\_cubic\_jordan22\_unipotent\_vertex.py}。
元26核、全rank1最高block、量子parameterの必要五次minor、三つのunit source表示、
最高PBW関係・全三次 $N$ 核・twist後の全point漸化式を検証する。
二次標準形・全被覆・PBWとtails・Serreの通常証明を有限標本で代替したものではない。
全CY包絡の中央三次切断とintrinsic幾何学的群oid・成分・相対降下は未完成である。
